% Modernised reading — fonds Alexandre Grothendieck, shelfmark 139, pages 1-2.
%
% This is an INTERPRETATION, not a transcription. In any dispute the
% transcription governs: whoever cites Grothendieck must cite batch-01.fr.tex.
%
% Pass: Opus 5.5 (claude-opus-5-5), 2026-10-10 — first pass, unchecked against
% the pages by a human. It was made on Opus 5.5 and has not been measured
% against the reference reading of folder 115 (made on Opus 5). The folder was
% read whole; it holds one batch, pages 1-2.
%
% This folder holds no manuscript and no mathematics: its two pages are
% photographs of one object, the Émile Picard medal of the Académie des
% sciences, struck in Grothendieck's name for 1977. The reading says so at the
% outset rather than manufacturing content; it has no spine section, and the
% four failure modes of the skill have nothing to bite on.
%
% Departures and choices, each footnoted:
%  - the identification of the portrait as Émile Picard comes from the reverse,
%    not from the obverse, which names no one (as the transcription notes);
%  - the transcription's « quatre lignes » are five once the rule is counted
%    separately, or four of text; the reading counts four lines of text;
%  - the accents the transcription restores (MÉDAILLE) are not on the metal;
%  - what is said of Picard, of the prize and of the engraver Morlon is cited
%    from memory and was not checked against the sources;
%  - the inventory files this folder in the group « Complexe de De Rham »
%    (136-139); nothing on the medal connects it to that subject, and the
%    reading does not invent a connection.
\documentclass[11pt,a4paper]{article}
\input{../preamble/grothendieck}

\folder{139}
\pages{1}{2}
\dating{1977}
\watermark{Édition de démonstration\\interprétation personnelle de l'œuvre}
\foldertitle{[Prix Émile Picard de l'Académie des Sciences de l'Institut de France] : médaille — lecture modernisée du dossier entier}

\begin{document}

\section*{Résumé}

\begin{resume}
Ce dossier ne contient aucune mathématique, ni même aucune écriture de
Grothendieck. Ses deux pages sont les deux faces photographiées d'une médaille :
celle du prix Émile Picard que l'Académie des sciences lui a décerné pour 1977.
L'avers porte un portrait de profil, signé du graveur ; le revers porte le nom
de l'Académie, celui de la médaille, celui du lauréat et le millésime.

Il n'y a donc rien à « moderniser » au sens où cette édition l'entend
d'ordinaire : pas d'énoncé à corriger, pas de notation à remettre au goût du
jour. Ce que cette lecture peut faire honnêtement, c'est dire ce qu'est
l'objet, ce qu'il porte, et où il se situe.

Il se situe à un moment singulier. En 1977, Grothendieck a quitté depuis
plusieurs années l'Institut des hautes études scientifiques et le monde
mathématique parisien ; il enseigne à Montpellier, loin des séminaires qui
avaient fait de lui la figure centrale de la géométrie algébrique. L'institution
le distingue alors qu'il s'en est éloigné. Qu'il ait conservé la médaille,
ou du moins qu'elle se soit retrouvée parmi ses papiers, est tout ce que le
dossier atteste : il ne dit rien de ce qu'il en pensait.

L'inventaire de Montpellier range ce dossier dans le groupe intitulé
« Complexe de De Rham », avec les dossiers 136 à 138. Rien sur la médaille ne
la rattache à ce sujet ; c'est la proximité de date, plutôt que de contenu, qui
semble l'avoir placée là.

\keywords{Émile Picard Medal, Académie des sciences, mathematical prizes}
\end{resume}

\section*{L'avers (page 1)}
\pagerange{1}{1}

Un portrait d'homme, de profil, tourné vers la gauche, en buste et en habit.
Aucune légende ne court autour. La seule inscription est la signature du
graveur, dans le champ, derrière l'épaule :
\begin{quote}
MORLON
\end{quote}
L'avers ne nomme pas le modèle ; c'est le revers, en titrant la pièce
« médaille Émile Picard », qui permet d'y reconnaître Picard.\footnote{La
transcription fait la même remarque : l'identification du portrait vient du
revers, la face elle-même ne la donne pas.} Émile Picard (1856--1941),
analyste --- le théorème de Picard sur les valeurs omises par une fonction
entière porte son nom --- fut secrétaire perpétuel de l'Académie des
sciences.\footnote{Cité de mémoire, comme tout ce qui est dit ici de Picard,
du prix et du graveur ; rien n'en a été vérifié pour cette lecture. La
signature MORLON est vraisemblablement celle du graveur en médailles
Pierre-Alexandre Morlon, actif dans la première moitié du vingtième
siècle, mais la pièce ne porte que le nom.}

\section*{Le revers (page 2)}
\pagerange{2}{2}

En légende circulaire, lue depuis le haut dans le sens des aiguilles d'une
montre :
\begin{quote}
ACADÉMIE DES SCIENCES DE L'INSTITUT DE FRANCE
\end{quote}
Au centre, quatre lignes de texte, les deux premières séparées des deux
dernières par un filet horizontal :\footnote{La transcription annonce
« quatre lignes que sépare un filet » et liste le filet comme un cinquième
élément ; on compte ici les lignes de texte.}
\begin{quote}
MÉDAILLE\footnote{Sans accent sur le métal, qui porte MEDAILLE et EMILE
PICARD en capitales sans diacritiques ; la transcription restitue l'accent
sur le premier mot seulement, et on la suit.}\\
EMILE PICARD\\
ALEXANDRE GROTHENDIECK\\
1977
\end{quote}
Une petite étoile est frappée en bas du champ, sous le millésime, là où
s'interrompt la légende circulaire.

La disposition dit à elle seule ce qu'est l'objet : en haut, l'institution et
la distinction, toujours les mêmes ; sous le filet, le nom du lauréat et
l'année, frappés pour lui. C'est un exemplaire nominatif d'une médaille dont
le coin principal sert d'un lauréat à l'autre.\footnote{Inférence tirée de la
seule mise en page de la pièce ; le dossier ne dit rien de la fabrication.}

\end{document}
